% GENERATED FILE. Edit canonical lessons, then run npm run generate:books.
% canonical-source-hash: fnv1a64:1cee1c7d5f38110f
% canonical-lessons: HI-C169-jorna, HI-C169-ghatana, HI-C169-barhana, HI-C169-bacana, HI-C169-bacna

\chapter{To join, To reduce, To increase, To save, To escape}
\label{ch:hi-to-join-to-reduce-to-increase-to-save-to-escape}
\hlchaptermodality{169}

\begin{chapteropening}
\textbf{By the end of this chapter:} \emph{I can say to join, to reduce, to increase, to save and to escape.}

Complete the last lesson of chapter 169: I can say to join, to reduce, to increase, to save and to escape.
\end{chapteropening}

\section[\texorpdfstring{\hi{जोड़ना} (joṛnā)}{जोड़ना (joṛnā)}]{\hi{जोड़ना} (joṛnā) — to join, to add}

\label{lesson:HI-C169-jorna}



\begin{quote}
Before the new one: say the Hindi for to drive, then the Hindi for to leave.
\end{quote}

\subsection*{You'll want to know: \hi{जोड़ना}}
\textbf{\hi{जोड़ना}} — \emph{joṛnā} — \textquotedblleft{}to join, to add\textquotedblright{}.

\begin{cousinweb}
\textbf{\hi{जोड़ना}} (\emph{joṛnā}) is a tadbhava, a word inherited and reshaped by sound change. It comes from \emph{*yuṭati}, \textquotedblleft{}is joined\textquotedblright{}, a form CDIAL reconstructs from its descendants (CDIAL 10496). CDIAL calls its link with Sanskrit \emph{yuj-}, \textquotedblleft{}to join\textquotedblright{}, obscure.
\end{cousinweb}

\begin{grammarlens}[title={What you've built}]
One more word for this chapter.
\end{grammarlens}

\subsection*{Your turn}
\begin{itemize}
\raggedright
  \item \emph{Say it:} \emph{joṛnā}
  \item \emph{Say it:} \emph{joṛnā}, once more
  \item \emph{Say it:} say \emph{joṛnā}
  \item \emph{Recall:} say the Hindi for to drive, then the Hindi for to leave, then say \emph{joṛnā} again
\end{itemize}

\begin{tcolorbox}[breakable,colback=teal!4,colframe=teal!35!black,title={Before you move on}]
What does \hi{जोड़ना} mean? (\textquotedblleft{}To join, to add\textquotedblright{}.) Say it once more.
\end{tcolorbox}
\section[\texorpdfstring{\hi{घटाना} (ghaṭānā)}{घटाना (ghaṭānā)}]{\hi{घटाना} (ghaṭānā) — to reduce, to subtract}

\label{lesson:HI-C169-ghatana}



\begin{quote}
Before the new one: say the Hindi for to leave, then the Hindi for to join.
\end{quote}

\subsection*{You'll want to know: \hi{घटाना}}
\textbf{\hi{घटाना}} — \emph{ghaṭānā} — \textquotedblleft{}to reduce, to subtract\textquotedblright{}.

\begin{cousinweb}
\textbf{\hi{घटाना}} (\emph{ghaṭānā}) is the causing form of \emph{ghaṭnā}, \textquotedblleft{}to decrease\textquotedblright{}, which CDIAL lists under a reconstructed \emph{*ghaṭṭati}, \textquotedblleft{}decreases, is wanting\textquotedblright{} (CDIAL 4415).
\end{cousinweb}

\begin{grammarlens}[title={What you've built}]
One more word for this chapter.
\end{grammarlens}

\subsection*{Your turn}
\begin{itemize}
\raggedright
  \item \emph{Say it:} \emph{ghaṭānā}
  \item \emph{Say it:} \emph{ghaṭānā}, once more
  \item \emph{Say it:} say \emph{ghaṭānā}
  \item \emph{Recall:} say the Hindi for to leave, then the Hindi for to join, then say \emph{ghaṭānā} again
\end{itemize}

\begin{tcolorbox}[breakable,colback=teal!4,colframe=teal!35!black,title={Before you move on}]
What does \hi{घटाना} mean? (\textquotedblleft{}To reduce, to subtract\textquotedblright{}.) Say it once more.
\end{tcolorbox}
\section[\texorpdfstring{\hi{बढ़ाना} (baṛhānā)}{बढ़ाना (baṛhānā)}]{\hi{बढ़ाना} (baṛhānā) — to increase, to raise}

\label{lesson:HI-C169-barhana}



\begin{quote}
Before the new one: say the Hindi for to join, then the Hindi for to reduce.
\end{quote}

\subsection*{You'll want to know: \hi{बढ़ाना}}
\textbf{\hi{बढ़ाना}} — \emph{baṛhānā} — \textquotedblleft{}to increase, to raise\textquotedblright{}.

\begin{cousinweb}
\textbf{\hi{बढ़ाना}} (\emph{baṛhānā}) is the causing form of \emph{baṛhnā}, \textquotedblleft{}to increase, grow\textquotedblright{}, which comes from Sanskrit \emph{vardhate}, \textquotedblleft{}grows, increases\textquotedblright{} (CDIAL 11376).
\end{cousinweb}

\begin{grammarlens}[title={What you've built}]
One more word for this chapter.
\end{grammarlens}

\subsection*{Your turn}
\begin{itemize}
\raggedright
  \item \emph{Say it:} \emph{baṛhānā}
  \item \emph{Say it:} \emph{baṛhānā}, once more
  \item \emph{Say it:} say \emph{baṛhānā}
  \item \emph{Recall:} say the Hindi for to join, then the Hindi for to reduce, then say \emph{baṛhānā} again
\end{itemize}

\begin{tcolorbox}[breakable,colback=teal!4,colframe=teal!35!black,title={Before you move on}]
What does \hi{बढ़ाना} mean? (\textquotedblleft{}To increase, to raise\textquotedblright{}.) Say it once more.
\end{tcolorbox}
\section[\texorpdfstring{\hi{बचाना} (bacānā)}{बचाना (bacānā)}]{\hi{बचाना} (bacānā) — to save, to rescue}

\label{lesson:HI-C169-bacana}



\begin{quote}
Before the new one: say the Hindi for to reduce, then the Hindi for to increase.
\end{quote}

\subsection*{You'll want to know: \hi{बचाना}}
\textbf{\hi{बचाना}} — \emph{bacānā} — \textquotedblleft{}to save, to rescue\textquotedblright{}.

\begin{cousinweb}
\textbf{\hi{बचाना}} (\emph{bacānā}) is the causing form of \emph{bacnā}, \textquotedblleft{}to escape\textquotedblright{}: to save is to make someone escape. CDIAL lists both under Sanskrit \emph{vañcati}, \textquotedblleft{}moves\textquotedblright{}, whose causing form meant \textquotedblleft{}escapes from\textquotedblright{} (CDIAL 11208).
\end{cousinweb}

\begin{grammarlens}[title={What you've built}]
One more word for this chapter.
\end{grammarlens}

\subsection*{Your turn}
\begin{itemize}
\raggedright
  \item \emph{Say it:} \emph{bacānā}
  \item \emph{Say it:} \emph{bacānā}, once more
  \item \emph{Say it:} say \emph{bacānā}
  \item \emph{Recall:} say the Hindi for to reduce, then the Hindi for to increase, then say \emph{bacānā} again
\end{itemize}

\begin{tcolorbox}[breakable,colback=teal!4,colframe=teal!35!black,title={Before you move on}]
What does \hi{बचाना} mean? (\textquotedblleft{}To save, to rescue\textquotedblright{}.) Say it once more.
\end{tcolorbox}
\section[\texorpdfstring{\hi{बचना} (bacnā)}{बचना (bacnā)}]{\hi{बचना} (bacnā) — to escape, to be saved}

\label{lesson:HI-C169-bacna}



\begin{quote}
Before the new one: say the Hindi for to increase, then the Hindi for to save.
\end{quote}

\subsection*{You'll want to know: \hi{बचना}}
\textbf{\hi{बचना}} — \emph{bacnā} — \textquotedblleft{}to escape, to be saved\textquotedblright{}.

\begin{cousinweb}
\textbf{\hi{बचना}} (\emph{bacnā}) is a tadbhava, a word inherited from Sanskrit and reshaped by sound change. It comes from \emph{vañcati}, \textquotedblleft{}moves; escapes from\textquotedblright{} (CDIAL 11208).
\end{cousinweb}

\begin{grammarlens}[title={What you've built}]
One more word for this chapter.
\end{grammarlens}

\subsection*{Your turn}
\begin{itemize}
\raggedright
  \item \emph{Say it:} \emph{bacnā}
  \item \emph{Say it:} \emph{bacnā}, once more
  \item \emph{Say it:} say \emph{bacnā}
  \item \emph{Recall:} say the Hindi for to increase, then the Hindi for to save, then say \emph{bacnā} again
\end{itemize}

\begin{tcolorbox}[breakable,colback=teal!4,colframe=teal!35!black,title={Before you move on}]
What does \hi{बचना} mean? (\textquotedblleft{}To escape, to be saved\textquotedblright{}.) Say it once more.
\end{tcolorbox}
